% sec4.tex — c08｜原文 p26–30｜§4 大 $N$ 极限下的指标、§4.1 黑洞、Cardy 极限、致谢
% 编号公式 (4.1)–(4.16) 共 16 条；编号环境：习题 4.1；无表格。
% [原文 p.26]
在本节中，我们将计算 $\sN=4$ $\SU(N)$ 超 Yang–Mills 理论在大 $N$ 极限下的超共形指标。其做法是用鞍点近似该积分。同样的鞍点技巧也适用于其他超对称规范理论的大 $N$ 极限 [21,25,40,41]，但这里我们只关注 $\sN=4$ 的情形。当分解为 $\sN=2$ 多重态时，$\sN=4$ 超 Yang–Mills 理论包含一个超多重态和一个矢量多重态，二者都在伴随表示中。单个 $\sN=1$ 矢量多重态的单字母指标为
\[
\begin{aligned}
i(p,q,x,s)&=\frac{-p-q+2pq+(pq)^{\frac13}x^{-1}-(pq)^{\frac23}x}{(1-p)(1-q)}\\
&\quad+\frac{(pq)^{\frac13}x^{\frac12}s-(pq)^{\frac23}x^{-\frac12}s^{-1}+(pq)^{\frac13}x^{\frac12}s^{-1}-(pq)^{\frac23}x^{-\frac12}s}{(1-p)(1-q)}\\
&=\frac{(pq)^{\frac13}(s\sqrt x+\frac{\sqrt x}{s}+\frac1x)-(pq)^{\frac23}(\frac{1}{s\sqrt x}+\frac{s}{\sqrt x}+x)-p-q+2pq}{(1-p)(1-q)} .
\end{aligned}
\tag{4.1}
\]
这里我们引入了逸度 $s$，它对应与 $\sN=4$ 的 $\sN=2$ 子代数对易的荷。它使两个半超多重态带电荷 $\pm1$。要得到 $\SU(N)$ $\sN=4$ 矢量多重态的指标，我们需要把这个指标乘以 $\chi_{\adj}(a_i)$，其中 $a_i$ 是与 $\SU(N)$ 对称性相伴的逸度。有趣的是，单个 $\sN=4$ 矢量的指标满足
\[
1-i(p,q,x,s)=\frac{(1-(pq)^{\frac13}s\sqrt x)(1-(pq)^{\frac13}\frac{\sqrt x}{s})(1-(pq)^{\frac13}\frac1x)}{(1-p)(1-q)} .
\tag{4.2}
\]
这将在后面用到。

规范理论的指标为
\[
\begin{aligned}
\sI&=\frac{1}{N!}\oint\Big(\prod_{j=1}^{N}\frac{\dd a_j}{2\pi\ii}\Big)\Delta(a_i)\,PE\big[i(p,q,x,s)\chi_{\adj}(a_i)\big]\\
&=\frac{1}{N!}\oint\Big(\prod_{j=1}^{N}\frac{\dd a_j}{2\pi\ii}\Big)\exp\Big(-\sum_{n=1}^{\infty}\frac1n\big(1-i(p^n,q^n,x^n,s^n)\big)\sum_{i\neq j}\frac{a_i^n}{a_j^n}\Big) .
\end{aligned}
\tag{4.3}
\]
这类矩阵模型在大 $N$ 极限下已在 [42] 以及 [12] 中被分析。我们遵循他们的分析。引入本征值密度 $\rho(\theta)$ 较为方便，其归一化为 $\int_{-\pi}^{\pi}\dd\theta\,\rho(\theta)=1$。指标于是成为关于 $\rho(\theta)$ 的泛函积分，但要求 $\rho$ 非负。该泛函积分的作用量为
\[
\begin{aligned}
\sI&=\int[\mathcal{D}\rho(\theta)]\exp\Big(N^2\int\dd\theta_1\dd\theta_2\,\rho(\theta_1)V(\theta_1-\theta_2)\rho(\theta_2)\Big)\\
V(\theta)&=-\sum_{n=1}^{\infty}\frac1n\big(1-i(p^n,q^n,x^n,s^n)\cos(n\theta)\big).
\end{aligned}
\tag{4.4}
\]
% [原文 p.27]
这些变量的好处在于，由于作用量中的 $N^2$ 因子，在大 $N$ 时该积分变成鞍点积分。我们进一步用密度模 $\rho_n=\int\dd\theta\,\rho(\theta)e^{in\theta}$ 与势 $V_n=\int\dd\theta\,V(\theta)\cos(n\theta)$ 来简化它。注意 $\rho_{-n}=\rho_n^\dagger$。泛函积分变为
\[
\begin{aligned}
\sI&=\int[\dd\rho_n\,\dd\rho_n^\dagger]\exp\Big(\frac{N^2}{2\pi}\sum_{n=1}^{\infty}\frac1n V_n|\rho_n|^2\Big)\\
V_n&=\frac{2\pi}{n}\big(1-i(p^n,q^n,x^n,s^n)\big).
\end{aligned}
\tag{4.5}
\]
由 (4.2) 可见，对所有 $n\ge1$ 都有 $V_n<0$。这使得所有模 $\rho_{n\ge1}$ 都获得质量。经典位形由 $\rho_0$ 主导。零模的作用量简单地为 $0$，积分由单圈行列式给出，
\[
\begin{aligned}
\sI&=\prod_{n=1}^{\infty}\frac{1}{1-i(p^n,q^n,x^n,s^n)}\\
&=\prod_{n=1}^{\infty}\frac{(1-p^n)(1-q^n)}{(1-(pq)^{\frac n3}s^n x^{\frac n2})(1-(pq)^{\frac n3}s^{-n}x^{\frac n2})(1-(pq)^{\frac n3}x^{-n})} .
\end{aligned}
\tag{4.6}
\]
容易看出，这个闭合形式的指标实际上是某个相对简单的函数的齐次指数。
\[
\begin{aligned}
\sI&=PE[i_G(p,q,x,s)]\\
i_G(p,q,x,s)&=\frac{(pq)^{\frac13}s\sqrt x}{1-(pq)^{\frac13}s\sqrt x}+\frac{(pq)^{\frac13}\frac{\sqrt x}{s}}{1-(pq)^{\frac13}\frac{\sqrt x}{s}}+\frac{(pq)^{\frac13}\frac1x}{1-(pq)^{\frac13}\frac1x}-\frac{p}{1-p}-\frac{q}{1-q} .
\end{aligned}
\tag{4.7}
\]
在大 $N$ 下，$\SU(N)$ $\sN=4$ 超 Yang–Mills 对偶于 $AdS_5\times S^5$ 上的弱耦合 IIB 型弦论。在对偶理论中，只有超引力态受保护（这是因为规范理论变强耦合时，弦模变得非常重）。这意味着 $\sN=4$ SYM 的指标应当计算 $AdS_5\times S^5$ 中超引力场的指标。指标 $\sI$ 的齐次对数，即 $i_G(p,q,x,s)$，确实就是 $AdS_5\times S^5$ 中超引力的单粒子指标 [12]。

\subsection{黑洞呢？}
$AdS_5\times S^5$ 中的超引力也允许超对称黑洞。它们具有 $\mathcal{O}(N^2)$ 的熵。由于我们发现大 $N$ 极限下指标是 $\mathcal{O}(1)$ 的，超对称黑洞的大量态在某种程度上对指标不可见。原因在于，黑洞的玻色态与费米态以这样的方式出现，以致在极大量的态之间发生严重抵消。如果能够避免这些抵消，我们就会找到一个作用量为 $\mathcal{O}(N^2)$ 的鞍点，而不是 $0$。这引出一个自然的问题：指标究竟会不会告诉我们黑洞态的信息？

这个问题的答案取决于玻色态与费米态之间抵消的性质。我们想强调的两种抵消，最好用一个例子
% [原文 p.28]
来说明。考虑指标的单变量极限 $p=q=t^3$、$x=s=1$。在此极限下单字母指标为
\[
i(t)=1-\frac{(1-t^2)^3}{(1-t^3)^2} .
\tag{4.8}
\]
让我们考虑某个大 $N$ 的 $\SU(N)$ 理论指标的 Taylor 级数展开
\[
\sI(t)=\sum_n a_n t^n .
\tag{4.9}
\]
玻色–费米态之间的严重抵消可能极其精细，使得该级数中的每个系数都是 $\mathcal{O}(1)$ 的小量而不是 $\mathcal{O}(N^2)$。在这种情况下，指标将永远无法探测与黑洞熵相伴的大量态简并。但另一种可能是抵消是粗糙的，即级数的系数是 $\mathcal{O}(N^2)$ 的大量，但符号快速振荡。即使在这种情况下，大的简并也难以捕捉，因为指标 $\sI(t)$ 只会看到平均掉的 $\mathcal{O}(1)$ 简并。然而，有一个简单的技巧可以绕开这个问题。若取 $t\to-t$（$t\in\mathbb{R}$），这或许能部分避开玻色–费米抵消，指标本身将是 $\mathcal{O}(N^2)$ 的。若 $a_n$ 的符号振荡周期为 2，则 $t\to-t$ 将完美避开所有抵消。若周期为 3，则我们可能需要取 $t\to t e^{2\pi i/3}$（$t\in\mathbb{R}$）以更好地避开抵消，如此等等。这一讨论的教训是：若玻色–费米态之间的抵消是粗糙的，那么也许可以通过为逸度打开非平庸相位来捕捉 $\mathcal{O}(N^2)$ 的简并增长。

最近，这一可能性已在 [43] 中、在所谓的指标「Cardy 极限」中得到了探索。对该分析而言，用下面描述的更「对称」的方式重新表述 $\sN=4$ 指标更为方便。$\sN=4$ 超共形代数的对易玻色生成元包括三个 R 对称生成元 $r_{i=1,2,3}$、两个转动生成元 $h_{i=1,2}$ 以及一个共形维数 $\Delta$。超荷用其在 Cartan 生成元下的荷 $Q^{r_1,r_2,r_3}_{h_1,h_2}$ 标号。这些荷取值 $\pm\frac12$，使得所有符号的乘积为 $+$。关于 $\mathcal{Q}^{+++}_{--}$ 的超共形指标为
\[
\sI=\Tr\, e^{-\sum_{i=1}^{3}\rho_i r_i}e^{-\sum_{i=1}^{2}\omega_i h_i},\qquad
\rho_1+\rho_2+\rho_3-\omega_1-\omega_2=2\pi i .
\tag{4.10}
\]
\begin{exercise}[4.1]
注意在式 (4.10) 中我们没有包含因子 $(-1)^F$。证明式 (4.10) 仍然定义了一个超共形指标。
提示：所有化学势的周期为 $4\pi i$。
\end{exercise}
% [原文 p.29]
$\sN=4$ 矢量多重态的单字母指标为
\[
i=1-2\,\frac{\sinh\frac{\rho_1}{2}\cdot\sinh\frac{\rho_2}{2}\cdot\sinh\frac{\rho_3}{2}}{\sinh\frac{\omega_1}{2}\cdot\sinh\frac{\omega_2}{2}} .
\tag{4.11}
\]
用这些变量，$\SU(N)$ 规范理论的指标 (4.3) 变为
\[
\sI=\frac{1}{N!}\int_0^{2\pi}\prod_{j=1}^{N}\dd\alpha_j\exp\Big(-\sum_{n=1}^{\infty}\frac2n\frac{\sinh\frac{n\rho_1}{2}\cdot\sinh\frac{n\rho_2}{2}\cdot\sinh\frac{n\rho_3}{2}}{\sinh\frac{n\omega_1}{2}\cdot\sinh\frac{n\omega_2}{2}}\sum_{i\neq j}e^{in\alpha_{ij}}\Big) .
\tag{4.12}
\]
我们不取大 $N$ 极限，而取「Cardy 极限」，即 $\operatorname{Re}(\omega_i)\to0^+$、$\operatorname{Im}(\omega_i)=0$。它被称为 Cardy 极限，是因为它让人联想到取高温极限，只不过在这里这个温度耦合的是角动量而不是能量。注意在此极限下 $\operatorname{Re}(\rho_i)\to0^+$，但由于约束 $\rho_1+\rho_2+\rho_3=\omega_1+\omega_2+2\pi i$ 以及对所有 $i=1,2,3$ 有 $\operatorname{Im}(\rho_i)=\mathcal{O}(1)$，并不需要 $\operatorname{Re}(\rho_i)=0$。在此极限下我们作替换 $\sinh(n\omega_i/2)\to n\omega_i/2$。式 (4.12) 变为
\[
\sI=\frac{1}{N!}\int_0^{2\pi}\prod_{j=1}^{N}\dd\alpha_j\exp\Big(-\frac{1}{\omega_1\omega_2}\sum_{n=1}^{\infty}\frac{2^3}{n^3}\Big(\sinh\frac{n\rho_1}{2}\cdot\sinh\frac{n\rho_2}{2}\cdot\sinh\frac{n\rho_3}{2}\Big)\sum_{i\neq j}e^{in\alpha_{ij}}\Big) .
\tag{4.13}
\]
指数可以进一步用三重对数 $\Li_3(x)\equiv\sum_{n=1}^{\infty}x^n/n^3$ 来化简。
\[
\begin{aligned}
&\sum_{n=1}^{\infty}\frac{2^3}{n^3}\Big(\sinh\frac{n\rho_1}{2}\cdot\sinh\frac{n\rho_2}{2}\cdot\sinh\frac{n\rho_3}{2}\Big)\sum_{i\neq j}e^{in\alpha_{ij}}
=\sum_{s_1,s_2,s_3=\pm1}\sum_{i\neq j}\Li_3\big(e^{\frac{s_i\rho_i}{2}+i\alpha_{ij}}\big)\\
&=\sum_{s_1s_2s_3=1}\sum_{i\neq j}\Big(\Li_3\big(e^{\frac{s_i\rho_i}{2}+i\alpha_{ij}}\big)-\Li_3\big(e^{-\frac{s_i\rho_i}{2}-i\alpha_{ij}}\big)\Big)
\end{aligned}
\tag{4.14}
\]
三重对数具有如下性质
\[
\Li_3(e^x)-\Li_3(e^{-x})=-\frac{x^3}{6}+\frac{\pi i x^2}{2}+\frac{\pi^2 x}{3},\qquad (0<\operatorname{Im}(x)<2\pi,\ \operatorname{Re}(x)\ge0) .
\tag{4.15}
\]
对 $2\pi p<\operatorname{Im}(x)<2\pi(p+1)$，我们需要在右端作平移 $x\to2\pi ip$。经过一番整理，可以注意到作用量泛函的鞍点在 $\alpha_i=\alpha$ 处。该鞍点处作用量的经典值按如下方式标度
\[
\log(\sI)\sim N^2\frac{\rho_1\rho_2\rho_3}{\omega_1\omega_2} .
\tag{4.16}
\]
这正是在 [44] 中被发现、用于产生 $AdS_5$ 中带电 BPS 黑洞正确熵的泛函。同一问题在 [45] 中通过把 Bethe ansatz 方法应用于超共形指标而得到另一种处理。关于有限 $N$ 下指标的数值研究，见 [46,47]。
% [原文 p.30]
\paragraph{致谢}
我们要感谢 DESY 理论物理组，特别是 Elli Pomoni 的盛情款待。我们还要感谢印度人民对基础科学研究的持续支持。

% [Tutor Note: 物理：$\sN=4$ SYM 的指标在大 $N$ 下计数的是 $AdS_5\times S^5$ 超引力态，因为弦模随耦合变强而变得非常重。]
% [Tutor Note: 易错：$s$ 是与 $\sN=2$ 子代数对易的那个 $\U(1)$ 的逸度；两个半超多重态带电荷 $\pm1$，故 (4.1) 中 $s$ 与 $s^{-1}$ 成对出现。]
% [Tutor Note: 易错：(4.1) 的 $x^{\frac12}$ 项对应 $r=1$ 的手征多重态、$x^{-\frac12}$ 项对应 $r=-1$，合并后化成三组因子。]
% [Tutor Note: 物理：(4.2) 的 $1-i$ 恰是 (4.6) 分子的来源；分母的三个因子分别对应 $x^{\frac12}s$、$x^{\frac12}/s$、$x^{-1}$ 的模。]
% [Tutor Note: 易错：(4.3) 中 $\Delta(a_i)$ 与 $PE$ 展开给出 $\sum_{i\neq j}a_i^n/a_j^n$，且不要漏掉前因子 $1/N!$。]
% [Tutor Note: 物理：作用量正比于 $N^2$，大 $N$ 时积分由鞍点主导；$V_n<0$ 使 $\rho_{n\ge1}$ 获得质量，只有零模 $\rho_0$ 无质量。]
% [Tutor Note: 易错：(4.5) 的 $V_n$ 带 $2\pi/n$ 归一化；对 $\rho_n$ 作高斯积分给出 $1/(1-i)$，再取对数即得 $PE$。]
% [Tutor Note: 物理：(4.6) 的闭合形式表明大 $N$ 指标就是 $PE[i_G]$；$i_G$ 与超引力 Kaluza–Klein 谱的单粒子指标一致 [12]。]
% [Tutor Note: 物理：黑洞熵为 $\mathcal{O}(N^2)$ 而指标为 $\mathcal{O}(1)$——典型黑洞态在指标中成对抵消，这正是「指标看不到黑洞」的根源。]
% [Tutor Note: 易错：Cardy 极限不是高温极限；此处温度耦合的是角动量，条件是 $\operatorname{Re}(\omega_i)\to0^+$ 而非 $\beta\to0$。]
% [Tutor Note: 易错：(4.10) 虽无 $(-1)^F$ 仍是指标，因为化学势以 $4\pi i$ 为周期，绕一圈给出的符号可重新解释为 $(-1)^F$（习题 4.1）。]
% [Tutor Note: 易错：(4.12)→(4.13) 只作 $\sinh(n\omega_i/2)\to n\omega_i/2$，并保留 $1/(\omega_1\omega_2)$ 与 $2^3/n^3$。]
% [Tutor Note: 易错：(4.14) 第二行下标 $s_1s_2s_3=1$ 只取乘积为 $+$ 的符号组合，配对 $s\to-s$ 后出现 $\Li_3(x)-\Li_3(1/x)$。]
% [Tutor Note: 文献：(4.16) 与 [44] 的带电 BPS 黑洞熵一致；Bethe ansatz 方法见 [45]；有限 $N$ 数值研究见 [46,47]。]
% [Tutor Note: 易错：原文 (4.4) 第二行括号不平衡，且 $\cos(n\theta)$ 应在因子 $(1-i(\cdots))$ 之外，正确形式见 ERRATA 候选；译文按平衡括号照录。]
% [Tutor Note: 易错：原文 (4.5) 第二行缺右括号，译文补为平衡括号；已登记 ERRATA。]
